% !TeX program = xelatex
% Native Beamer reconstruction of presentation.pdf.
% Build: xelatex -interaction=nonstopmode -halt-on-error -jobname=presentation_tex presentation.tex
\documentclass[aspectratio=169,10pt]{beamer}

\usepackage{fontspec}
\usepackage{graphicx}
\usepackage{tikz}
\usepackage{xcolor}
\usepackage{array}
\usepackage{tabularx}

\setsansfont{Noto Sans CJK KR}
\setmainfont{Noto Sans CJK KR}
\setbeamertemplate{navigation symbols}{}
\setbeamertemplate{headline}{}
\setbeamertemplate{footline}{}
\setbeamersize{text margin left=0pt,text margin right=0pt}

\definecolor{Ink}{HTML}{191919}
\definecolor{Canvas}{HTML}{F2F2F2}
\definecolor{AccentGreen}{HTML}{27DB3E}
\definecolor{AccentRed}{HTML}{FF3838}
\definecolor{BrandAccent}{HTML}{414AFF}

\newcommand{\CandidateName}{장성원}
\newcommand{\TargetCompany}{지원 기업}
\newcommand{\TargetRole}{Robotics Vision Engineer}
% LaTeX PDF Preview builds only a cached copy of this .tex file. Use the source
% deck's absolute path so image crops also resolve inside that cache directory.
\newcommand{\SourceDeck}{/home/swlinux/Desktop/workspace/ai-job-search/cv/ppt/presentation.pdf}
\newcommand{\AssetDir}{/home/swlinux/Desktop/workspace/ai-job-search/cv/ppt/assets}

% Absolute-position helpers. Coordinates are fractions of the 16:9 page.
\newcommand{\place}[4]{%
  \node[anchor=north west,inner sep=0pt,text width=#3\paperwidth,align=left]
    at ([xshift=#1\paperwidth,yshift=-#2\paperheight]current page.north west) {#4};%
}
\newcommand{\cropimg}[6]{%
  \includegraphics[page=#1,trim=#2bp #3bp #4bp #5bp,clip,width=#6\paperwidth]{\SourceDeck}%
}
\newcommand{\canvas}{%
  \fill[Canvas] (current page.south west) rectangle (current page.north east);%
}
\newcommand{\deckbar}{%
  \fill[black] (current page.north west)
    rectangle ([yshift=-0.145\paperheight]current page.north east);%
}
\newcommand{\deckheader}[2][30]{%
  \deckbar
  \place{0.01}{0.018}{0.97}{\color{white}\bfseries\fontsize{#1}{#1}\selectfont #2}%
}
\newcommand{\arrowmark}{%
  \node[anchor=center,font=\fontsize{18}{18}\selectfont]
    at ([xshift=-0.065\paperwidth,yshift=0.11\paperheight]current page.south east) {$\rightarrow$};%
}
\newcommand{\bulletline}[1]{\textbullet\enspace #1\par}

% English Mobin pages use the existing Astra evidence captures.
\newcommand{\PhyAssetDir}{/home/swlinux/Desktop/workspace/ai-job-search/cv/ppt/assets/phy}
\newcommand{\flowstep}[5]{%
  \node[anchor=north west,inner sep=0pt,text width=#3\paperwidth,
    minimum height=0.14\paperheight,align=center,fill=white,draw=Ink,line width=0.65pt]
    at ([xshift=#1\paperwidth,yshift=-#2\paperheight]current page.north west)
    {\fontsize{7.2}{10}\selectfont\textbf{#4}\\[3pt]\fontsize{6}{8}\selectfont #5};%
}
\newcommand{\flowarrow}[3]{%
  \draw[->,line width=0.7pt] ([xshift=#1\paperwidth,yshift=-#3\paperheight]current page.north west)
    -- ([xshift=#2\paperwidth,yshift=-#3\paperheight]current page.north west);%
}

\begin{document}

% 1. Cover
\begin{frame}[plain]
\begin{tikzpicture}[remember picture,overlay]
  \fill[white] (current page.south west) rectangle (current page.north east);
  \fill[black] (current page.north west)
    rectangle ([yshift=-0.265\paperheight]current page.north east);
  \place{0.027}{0.035}{0.47}{\color{white}\bfseries\fontsize{70}{30}\selectfont\addfontfeatures{LetterSpace=-4}PUBLIC}
  \place{0.027}{0.270}{1.20}{\color{Ink}\bfseries\fontsize{70}{30}\selectfont\addfontfeatures{LetterSpace=-4}RELATIONS}
  \place{0.45}{0.835}{0.23}{\fontsize{10}{12}\selectfont Presented By :\\[-1pt]\mbox{\textbf{Seongwon Jang}}}
  \place{0.73}{0.835}{0.12}{\fontsize{10}{12}\selectfont Present To :}
  \arrowmark
\end{tikzpicture}
\end{frame}

% 2. About
\begin{frame}[plain]
\begin{tikzpicture}[remember picture,overlay]
  \canvas \deckheader{1. 자기 소개}
  \place{0.136}{0.235}{0.122}{\cropimg{2}{196}{400}{1068}{190}{0.122}}
  \place{0.145}{0.515}{0.12}{\centering\bfseries\fontsize{9}{11}\selectfont 장성원\\[-1pt]\mdseries 2000.11.23.}
  \draw[line width=1.1pt] ([xshift=0.385\paperwidth,yshift=-0.21\paperheight]current page.north west)
    -- ([xshift=0.385\paperwidth,yshift=-0.95\paperheight]current page.north west);
  \place{0.040}{0.610}{0.027}{\includegraphics[width=\linewidth]{\AssetDir/github-logo.png}}
  \place{0.040}{0.697}{0.027}{\includegraphics[width=\linewidth]{\AssetDir/huggingface-logo.png}}
  \place{0.040}{0.784}{0.027}{\includegraphics[width=\linewidth]{\AssetDir/dacon-logo.png}}
  \place{0.040}{0.871}{0.027}{\includegraphics[width=\linewidth]{\AssetDir/linkedin-logo.png}}
  \place{0.075}{0.61}{0.29}{\fontsize{6.6}{9}\selectfont
    \textbf{GitHub}\\http://github.com/JungSeong\par\vspace{4pt}
    \textbf{HuggingFace}\\http://huggingface.co/JungSeong2\par\vspace{4pt}
    \textbf{Dacon}\\http://dacon.io/myprofile/537952/home\par\vspace{4pt}
    \textbf{LinkedIN}\\http://www.linkedin.com/in/seongwonjang}
  \place{0.405}{0.225}{0.54}{\fontsize{7.3}{9.5}\selectfont
    \textbf{학력}\\[-1pt]
    \textbf{26.02} $\mid$ 세종대학교 학사 졸업\\
    평점 4.02/4.5 (전공 평점 4.13/4.5)\par\vspace{5pt}
    \textbf{자격/수상}\\[-1pt]
    \textbf{2025.12} $\mid$ 정보처리기사 \quad \textbf{2020.11} $\mid$ 자동차운전면허 2종 보통\\
    \textbf{2025.06} $\mid$ ROS2와 AI를 활용한 자율주행\&로봇팔 개발자 부트캠프 최우수상\\
    \textbf{2024.12} $\mid$ 2024 IN-JEJU Challenge 해커톤 우수상\par\vspace{5pt}
    \textbf{경력}\\[-1pt]
    \textbf{25.07--09} $\mid$ UNIST 연구 인턴십\\
    \bulletline{Mixture-of-Expert 모델의 효율적인 파라미터 파인튜닝 기법 연구 개발}
    \vspace{5pt}\textbf{주요 프로젝트}\\[-1pt]
    \bulletline{AI for Industry Challenge}
    \bulletline{RoboDine}
    \bulletline{구조물 안정성 물리 추론 AI 경진대회}
    \vspace{5pt}\textbf{학내외 활동}\\[-1pt]
    \textbf{2026} $\mid$ 모두의 연구소 Advanced Computer Vision Lab 스터디원\\
    \textbf{2025} $\mid$ 오로카 판교 로봇 스터디원\\
    \textbf{2024} $\mid$ 교내 프로그래밍 동아리 Prompt C언어 팀장}
  \arrowmark
\end{tikzpicture}
\end{frame}

% 3. Skills
\begin{frame}[plain]
\begin{tikzpicture}[remember picture,overlay]
  \canvas \deckheader{1. 자기 소개}
  \place{0.01}{0.155}{0.30}{\bfseries\fontsize{13}{18}\selectfont 기술 스택}
  \place{0.01}{0.25}{0.08}{\bfseries\fontsize{11}{15}\selectfont HW}
  \place{0.028}{0.315}{0.14}{\cropimg{3}{40}{323}{1203}{254}{0.14}}
  \place{0.2071}{0.315}{0.1258}{\cropimg{3}{287}{323}{976}{254}{0.1258}}
  \place{0.33225}{0.315}{0.1905}{\cropimg{3}{482}{323}{690}{254}{0.1905}}
  \place{0.008}{0.620}{0.18}{\centering\fontsize{5.8}{8}\selectfont\mbox{Universal Robotics (UR5e)}}
  \place{0.22}{0.620}{0.16}{\centering\fontsize{6.5}{8}\selectfont MyCobot280}
  \place{0.42}{0.620}{0.17}{\centering\fontsize{6.5}{8}\selectfont Pinky}
  \place{0.01}{0.72}{0.08}{\bfseries\fontsize{11}{15}\selectfont SW}
  \place{0.02}{0.775}{0.50}{\cropimg{3}{31}{71}{694}{610}{0.50}}
  \draw[line width=1.1pt] ([xshift=0.55\paperwidth,yshift=-0.17\paperheight]current page.north west)
    -- ([xshift=0.55\paperwidth,yshift=-0.95\paperheight]current page.north west);
  \place{0.562}{0.34}{0.42}{\fontsize{8}{11}\selectfont
    \textbf{교육}\\
    \textbf{2025.04--06} $\mid$ ROS2와 AI를 활용한 자율주행\&로봇팔 직업 훈련\par\vspace{8pt}
    \textbf{자격/면허}\\
    \textbf{2026.03} $\mid$ 정보처리기사\\
    \textbf{2020.10} $\mid$ 2종 보통 면허\par\vspace{8pt}
    \textbf{어학}\\
    \textbf{2026.03} $\mid$ 오픽 IH\\
    \textbf{2025.06} $\mid$ 토익 925}
\end{tikzpicture}
\end{frame}

% 4. Projects index
\begin{frame}[plain]
\begin{tikzpicture}[remember picture,overlay]
  \canvas \deckheader{2. 주요 프로젝트}
  \place{0.025}{0.42}{0.90}{\fontsize{16}{21}\selectfont
    2-1. AI for Industry Challenge}
  \place{0.025}{0.50}{0.90}{\fontsize{16}{21}\selectfont
    2-2. RoboDine}
  \place{0.025}{0.58}{0.90}{\fontsize{16}{21}\selectfont
    2-3. 구조물 안정성 물리추론 AI 경진대회}
\end{tikzpicture}
\end{frame}

% 5. AIC summary
\begin{frame}[plain]
\begin{tikzpicture}[remember picture,overlay]
  \canvas \deckheader{2-1. AI for Industry Challenge}
  \place{0.006}{0.33}{0.61}{\fontsize{7.5}{10}\selectfont
    \renewcommand{\arraystretch}{1.35}
    \begin{tabular}{|>{\centering\arraybackslash}m{0.075\paperwidth}|>{\raggedright\arraybackslash}m{0.477\paperwidth}|}\hline
      담당 역할 & 팀장 \\ \hline
      주요 업무 &
      \bulletline{Gazebo Simulator 기반 랜덤화 환경 및 정책별 데이터 수집 ROS2 노드 개발}
      \bulletline{\textbf{YOLO Keypoint Detection 및 Multiview Triangulation} 기반 포트 3-DoF 추론 로직 구현}
      \bulletline{Fine-Grained한 상황에서 케이블 정렬을 위한 \textbf{MultiView Bilinear Cross-Attention Regressor} 구현}
      \bulletline{팀원 구인 및 GitHub, HuggingFace, Notion 관리} \\ \hline
      기술 스택 &
      HW : Universal Robotics UR5e\par
      SW : Python, PyTorch, Ubuntu 24.04, ROS2 Kilted Kaiju, Gazebo, Docker, YOLOv8, OpenCV\par
      GPU : DGX Spark, NVIDIA RTX 3060 \\ \hline
    \end{tabular}}
  \draw[line width=1.1pt] ([xshift=0.625\paperwidth,yshift=-0.215\paperheight]current page.north west)
    -- ([xshift=0.625\paperwidth,yshift=-0.95\paperheight]current page.north west);
  \place{0.635}{0.351}{0.35}{\cropimg{5}{915}{141}{17}{249}{0.35}}
  \place{0.635}{0.888}{0.35}{\makebox[\linewidth][c]{\bfseries\fontsize{9}{11}\selectfont Peg-In-Hole}}
\end{tikzpicture}
\end{frame}

% 6. AIC randomization
\begin{frame}[plain]
\begin{tikzpicture}[remember picture,overlay]
  \canvas \deckheader{2-1. AI for Industry Challenge}
  \place{0.012}{0.175}{0.965}{\fontsize{7.5}{10}\selectfont
    1. Gazebo Simulator에서 동작하는 \textbf{도메인 랜덤화 데이터 세트 생성 및 수집 ROS2 노드 구현}\\[-1pt]
    \bulletline{Uniform Distribution으로 Task Board의 XY/yaw, 포트 타입 및 XY/yaw 변화, Gaussian Distribution으로 조명 변화를 둔 시나리오 생성과 HuggingFace Hub 업로드 로직 구현}}
  \draw[line width=1.1pt] ([xshift=0.46\paperwidth,yshift=-0.29\paperheight]current page.north west)
    -- ([xshift=0.46\paperwidth,yshift=-0.92\paperheight]current page.north west);
  \place{0.008}{0.445}{0.44}{\cropimg{6}{11}{110}{795}{362}{0.44}}
  \place{0.465}{0.34}{0.52}{\cropimg{6}{675}{111}{17}{278}{0.52}}
  \place{0.02}{0.895}{0.42}{\makebox[\linewidth][c]{\bfseries\fontsize{8.5}{10.5}\selectfont Distribution (Task Board, Port, Lightening)}}
  \place{0.47}{0.895}{0.51}{\makebox[\linewidth][c]{\bfseries\fontsize{8.5}{10.5}\selectfont Data Collector}}
\end{tikzpicture}
\end{frame}

% 7. AIC triangulation
\begin{frame}[plain]
\begin{tikzpicture}[remember picture,overlay]
  \canvas \deckheader{2-1. AI for Industry Challenge}
  \place{0.012}{0.165}{0.965}{\fontsize{7.5}{10}\selectfont
    2. \textbf{YOLO Keypoint Detection 및 Multiview Triangulation 기반 포트 3-DoF 추론 로직 구현}, 5개 Case에서 MAE 약 5mm 내외로 예측 정확도 달성}
  \draw[line width=1.1pt] ([xshift=0.58\paperwidth,yshift=-0.215\paperheight]current page.north west)
    -- ([xshift=0.58\paperwidth,yshift=-0.95\paperheight]current page.north west);
  \place{0.009}{0.245}{0.555}{\cropimg{7}{13}{120}{627}{198}{0.555}}
  \place{0.59}{0.25}{0.39}{\cropimg{7}{856}{40}{30}{201}{0.39}}
\end{tikzpicture}
\end{frame}

% 8. AIC architecture
\begin{frame}[plain]
\begin{tikzpicture}[remember picture,overlay]
  \canvas \deckbar
  \place{0.012}{0.165}{0.965}{\fontsize{8}{10.5}\selectfont
    3. Fine-Grained한 상황에서 케이블 정렬을 위한 \textbf{MultiView Bilinear Cross-Attention Regressor} 구현\\[-1pt]
    \bulletline{\textbf{Worst Case 대비 XYZ MAE 21.3\% 및 RPY MAE 16.5\% 향상}}}
  \place{0.006}{0.27}{0.985}{\cropimg{8}{8}{58}{20}{217}{0.985}}
\end{tikzpicture}
\end{frame}

% 9. AIC result
\begin{frame}[plain]
\begin{tikzpicture}[remember picture,overlay]
  \canvas \deckheader{2-1. AI for Industry Challenge}
  \place{0.012}{0.165}{0.965}{\fontsize{8}{10.5}\selectfont
    3. Fine-Grained한 상황에서 케이블 정렬을 위한 \textbf{MultiView Bilinear Cross-Attention Regressor} 구현\\[-1pt]
    \bulletline{\textbf{Case에서 Worst Case 대비 XYZ MAE 21.3\% 및 RPY MAE 16.5\% 향상}}}
  \draw[line width=1.1pt] ([xshift=0.655\paperwidth,yshift=-0.225\paperheight]current page.north west)
    -- ([xshift=0.655\paperwidth,yshift=-0.96\paperheight]current page.north west);
  \place{0.016}{0.35}{0.625}{\cropimg{9}{23}{135}{517}{282}{0.625}}
  \place{0.665}{0.305}{0.33}{\cropimg{9}{948}{80}{7}{250}{0.33}}
\end{tikzpicture}
\end{frame}

% 10. RoboDine summary
\begin{frame}[plain]
\begin{tikzpicture}[remember picture,overlay]
  \canvas \deckheader{2-2. RoboDine}
  \place{0.006}{0.375}{0.978}{\fontsize{8}{11}\selectfont
    \renewcommand{\arraystretch}{1.45}
    \begin{tabularx}{\linewidth}{|>{\centering\arraybackslash}p{0.075\paperwidth}|X|}\hline
      담당 역할 & 팀원 \\ \hline
      주요 업무 &
      \bulletline{데이터 전처리 및 파라미터 튜닝을 통한 단안 2D 카메라가 달린 로봇팔의 객체 3-DoF 추정 정확도 14\% 향상}
      \bulletline{단안 카메라 및 2D LiDAR 거리 정보를 융합한 ROS2 기반 장애물 인식 및 회피 노드 설계} \\ \hline
      기술 스택 & HW : MyCobot280, Pinky\\
      & SW : Python, PyTorch, Ubuntu 24.04, ROS2 Jazzy, OpenCV, YOLOv8\\
      & GPU : NVIDIA RTX 3060 \\ \hline
    \end{tabularx}}
\end{tikzpicture}
\end{frame}

% 11. RoboDine pose estimation
\begin{frame}[plain]
\begin{tikzpicture}[remember picture,overlay]
  \canvas \deckheader{2-2. RoboDine}
  \place{0.014}{0.165}{0.965}{\fontsize{7.5}{10}\selectfont
    1. \textbf{데이터 전처리 및 파라미터 튜닝을 통한 단안 2D 카메라가 달린 로봇팔의 객체 3-DoF 추론 정확도 14\% 향상}\\[-1pt]
    \bulletline{데이터 증강, Dropout, Pose Loss 계수 상향해 YOLO pose 모델 학습}
    \bulletline{카메라 캘리브레이션 내부 파라미터와 YOLO 추론값으로 solvePnP 연산}}
  \place{0.008}{0.315}{0.98}{\cropimg{11}{11}{35}{21}{285}{0.98}}
\end{tikzpicture}
\end{frame}

% 12. RoboDine obstacle avoidance
\begin{frame}[plain]
\begin{tikzpicture}[remember picture,overlay]
  \canvas \deckbar
  \place{0.01}{0.165}{0.53}{\fontsize{5.5}{7.5}\selectfont
    2. ROS2 기반 단안 카메라 및 2D LiDAR 거리 정보를 융합한 장애물 인식 및 회피 노드 설계\\[-1pt]
    \bulletline{\textbf{라이다 거리 30CM 이내 + 동적 장애물}: 정지 이후 장애물이 시야에서 사라지면 이동}
    \bulletline{\textbf{라이다 거리 20CM 이내 + 정적 장애물}: 회피 이후 이동}}
  \draw[line width=1.1pt] ([xshift=0.537\paperwidth,yshift=-0.275\paperheight]current page.north west)
    -- ([xshift=0.537\paperwidth,yshift=-0.97\paperheight]current page.north west);
  \place{0.01}{0.42}{0.515}{\cropimg{12}{14}{161}{684}{343}{0.515}}
  \place{0.55}{0.205}{0.435}{\cropimg{12}{793}{31}{24}{166}{0.435}}
\end{tikzpicture}
\end{frame}

% 13. Mobin overview
\begin{frame}[plain]
\begin{tikzpicture}[remember picture,overlay]
  \canvas \deckheader{Mobin}
  \place{0.025}{0.175}{0.95}{\bfseries\fontsize{13}{17}\selectfont From sensors to person tracking and driving decisions}
  \flowstep{0.025}{0.30}{0.205}{Camera + LiDAR}{Instance mask + point cloud}
  \flowstep{0.275}{0.30}{0.205}{3D person observation}{Point association + ground removal}
  \flowstep{0.525}{0.30}{0.205}{Person tracking}{Position + velocity + uncertainty}
  \flowstep{0.775}{0.30}{0.200}{Driving decision}{CPA + distance + observation age}
  \flowarrow{0.230}{0.275}{0.37}\flowarrow{0.480}{0.525}{0.37}\flowarrow{0.730}{0.775}{0.37}
  \place{0.03}{0.52}{0.47}{\includegraphics[width=\linewidth,height=0.35\paperheight,keepaspectratio]{\PhyAssetDir/mobin_association.png}}
  \place{0.54}{0.54}{0.43}{\fontsize{8.5}{11}\selectfont
    \textbf{Observation processing}\\[5pt]
    Use TF at the point-cloud timestamp; associate person points after ground removal.\\[8pt]
    \textbf{Tracking and decision}\\[5pt]
    Update odom-frame tracks with Hungarian assignment; pass position and velocity to the decision logic.}
  \place{0.025}{0.935}{0.95}{\color{gray}\fontsize{5.3}{7}\selectfont Existing capture: YOLO-only and distance-associated observations; not every detection has a valid distance.}
\end{tikzpicture}
\end{frame}

% 14. Mobin limitations
\begin{frame}[plain]
\begin{tikzpicture}[remember picture,overlay]
  \canvas \deckheader{Mobin}
  \place{0.025}{0.175}{0.95}{\bfseries\fontsize{13}{17}\selectfont Execution results and remaining observation and tracking limits}
  \place{0.025}{0.265}{0.95}{\fontsize{8}{11}\selectfont Existing 60-second seeded-multi-person-tracking run. This is a single-run result, not a full performance statistic.}
  \place{0.035}{0.36}{0.43}{\fontsize{8.5}{11}\selectfont
    \renewcommand{\arraystretch}{1.48}
    \begin{tabularx}{\linewidth}{|X|r|}\hline
      \textbf{Metric} & \textbf{Recorded value} \\ \hline
      Common sensor ROI coverage & 16.67\% \\ \hline
      IDF1 & 0.0595 \\ \hline
      ID switches & 2 \\ \hline
      Ground-truth misses (FN) & 130 \\ \hline
      Patrol state & not\_started \\ \hline
    \end{tabularx}}
  \place{0.53}{0.38}{0.44}{\includegraphics[width=\linewidth,height=0.33\paperheight,keepaspectratio]{\PhyAssetDir/mobin_stale.png}}
  \place{0.53}{0.735}{0.44}{\fontsize{8}{11}\selectfont
    \textbf{Interpretation}\\[4pt]
    Common observation coverage and tracking quality remain limited. The run also recorded safety behavior triggered by observation age.}
  \place{0.025}{0.955}{0.95}{\color{gray}\fontsize{5.3}{7}\selectfont Values are taken from summary.json. They do not establish completed root-cause analysis or safe autonomous driving.}
\end{tikzpicture}
\end{frame}

% 13. Physics reasoning summary
\begin{frame}[plain]
\begin{tikzpicture}[remember picture,overlay]
  \canvas \deckheader[21]{2-3. 구조물 안정성 물리 추론 AI 경진대회}
  \place{0.007}{0.385}{0.985}{\fontsize{7.5}{10}\selectfont
    \renewcommand{\arraystretch}{1.42}
    \begin{tabularx}{\linewidth}{|>{\centering\arraybackslash}p{0.075\paperwidth}|X|}\hline
      담당 역할 & 팀원 \\ \hline
      주요 업무 &
      \bulletline{Multi-View Bidirectional Cross Attention 모델 설계 및 적용해 기존 Feature Fusion 모델 대비 logloss 48.2\% 및 Accuracy 3.13\% 개선}
      \bulletline{VideoMAE로 영상 특징 벡터 추출 및 Distillation Regularization 추가해 Baseline 대비 logloss 86.3\% 및 accuracy 9\% 개선}
      \bulletline{Backbone, Epochs, Batch Size 등 HyperParameter Search 수행} \\ \hline
      기술 스택 & SW : Python, PyTorch, Ubuntu 24.04\\
      & GPU : DGX Spark \\ \hline
    \end{tabularx}}
\end{tikzpicture}
\end{frame}

% 14. Physics reasoning cross attention
\begin{frame}[plain]
\begin{tikzpicture}[remember picture,overlay]
  \canvas \deckbar
  \place{0.01}{0.16}{0.97}{\fontsize{8}{10.5}\selectfont
    1. \textbf{Multi-View Bidirectional Cross Attention} 모델 설계 및 적용\\[-1pt]
    \bulletline{기존 Feature Fusion 모델 대비 \textbf{logloss 48.2\% 및 Accuracy 3.13\% 개선}}}
  \place{0.045}{0.235}{0.94}{\cropimg{14}{65}{89}{19}{191}{0.94}}
\end{tikzpicture}
\end{frame}

% 15. Physics reasoning distillation/search
\begin{frame}[plain]
\begin{tikzpicture}[remember picture,overlay]
  \canvas \deckheader[21]{2-3. 구조물 안정성 물리 추론 AI 경진대회}
  \place{0.01}{0.16}{0.49}{\fontsize{6.6}{9}\selectfont
    2. VideoMAE로 영상에서 특징 벡터 추출 및 Distillation Regularization 추가해 학습\\[-1pt]
    \bulletline{Baseline 대비 \textbf{logloss 86.3\% 및 accuracy 9\% 개선}}}
  \place{0.542}{0.16}{0.45}{\fontsize{6.6}{9}\selectfont
    3. Backbone, Epochs, Batch Size 등 HyperParameter Search 수행\\[-1pt]
    \bulletline{\textbf{Worst Trial 대비 logloss 95.5\% 및 accuracy 40\% 개선}}
    ※ CosineAnnealingLR 스케줄러 사용, Epochs에 따라서도 수렴 속도가 달라짐}
  \draw[line width=1.1pt] ([xshift=0.525\paperwidth,yshift=-0.235\paperheight]current page.north west)
    -- ([xshift=0.525\paperwidth,yshift=-0.975\paperheight]current page.north west);
  \place{0.075}{0.27}{0.375}{\cropimg{15}{110}{20}{795}{255}{0.375}}
  \place{0.57}{0.31}{0.385}{\cropimg{15}{819}{20}{68}{285}{0.385}}
\end{tikzpicture}
\end{frame}

\end{document}
